\documentclass[letterpaper,journal]{IEEEtran}

\usepackage{amsmath}
\usepackage{amsfonts}

\hyphenation{op-tical net-works semi-conduc-tor}

\begin{document}

\title{Learnable Full-Dimensional Emotion Prototypes with Angular Relational Learning for Multimodal Emotion Recognition in Conversation}
\author{}
\maketitle

\begin{abstract}
\end{abstract}

\section{Introduction}
\label{sec:introduction}

\section{Related Work}
\label{sec:related_work}

\section{Problem Definition and Proposed Method}
\label{sec:problem_and_method}

\subsection{Problem Definition}
\label{subsec:problem_definition}

\subsubsection{Multimodal Emotion Recognition in Conversation}

Let a conversational dataset be denoted by
\(\mathcal{D}=\{\mathcal{C}_{n}\}_{n=1}^{N}\), where conversation
\(\mathcal{C}_{n}=\{u_{n,t}\}_{t=1}^{L_n}\) contains an ordered sequence of
utterances. Each utterance is associated with aligned textual, acoustic, and
visual observations,
\begin{equation}
    x_{n,t}=\left(x_{n,t}^{T},x_{n,t}^{A},x_{n,t}^{V}\right),
    \qquad
    y_{n,t}\in\mathcal{Y}=\{1,\ldots,C\},
    \label{eq:rgrpl_problem_input}
\end{equation}
where \(y_{n,t}\) is its categorical emotion label. The modality set is
\(\mathcal{M}=\{T,A,V\}\). The objective is to learn a contextual classifier
\(f_{\Theta}\) that predicts the emotion of every valid utterance by exploiting
both conversational context and complementary multimodal evidence:
\begin{equation}
    \hat y_{n,t}
    =\underset{c\in\mathcal{Y}}{\arg\max}\;
    p_{\Theta}\!\left(
    c\mid x_{n,1:L_n}^{T},x_{n,1:L_n}^{A},x_{n,1:L_n}^{V}\right).
    \label{eq:rgrpl_problem_prediction}
\end{equation}
Padding utterances introduced by mini-batching are excluded from every
utterance-level loss using a binary validity mask.

\subsubsection{SDT Backbone and Learning Objective}

For each valid utterance \(i\), the SDT backbone produces a pure self-modal
representation \(h_i^m\in\mathbb{R}^{d_h}\) and an
interaction-enhanced representation
\(\widetilde h_i^m\in\mathbb{R}^{d_h}\) for every
\(m\in\mathcal{M}\). A feature-wise SDT gate fuses the enhanced
representations:
\begin{align}
    a_i^m &= G\widetilde h_i^m,
    \label{eq:rgrpl_sdt_gate_score}\\
    [\beta_i^m]_k
    &=\frac{\exp([a_i^m]_k)}
    {\sum_{q\in\mathcal{M}}\exp([a_i^q]_k)},
    \label{eq:rgrpl_sdt_gate}\\
    \widetilde h_i
    &=\sum_{m\in\mathcal{M}}\beta_i^m\odot\widetilde h_i^m,
    \label{eq:rgrpl_sdt_fusion}\\
    p_i^{\mathrm{SDT}}
    &=\operatorname{softmax}(W_o\widetilde h_i+b_o),
    \label{eq:rgrpl_sdt_output}
\end{align}
where \(k\in\{1,\ldots,d_h\}\) indexes feature dimensions and \(\odot\)
denotes element-wise multiplication.

The SDT training objective combines the class-weighted loss of the fused
classifier, three unimodal student losses, and self-distillation from the
fused prediction to the modality students:
\begin{equation}
    \mathcal{L}_{\mathrm{SDT}}
    =\gamma_1\mathcal{L}_{\mathrm{task}}
    +\gamma_2\mathcal{L}_{\mathrm{student}}
    +\gamma_3\mathcal{L}_{\mathrm{KD}}.
    \label{eq:rgrpl_sdt_objective}
\end{equation}
For a set \(\mathcal{I}\) of valid utterances in a mini-batch and class
weights \(\omega_c\), the first two terms are
\begin{align}
    \mathcal{L}_{\mathrm{task}}
    &=-\frac{1}{Z_{\mathcal{I}}}
    \sum_{i\in\mathcal{I}}\omega_{y_i}
    \log p_i^{\mathrm{SDT}}[y_i],
    \label{eq:rgrpl_task_loss}\\
    \mathcal{L}_{\mathrm{student}}
    &=-\sum_{m\in\mathcal{M}}\frac{1}{Z_{\mathcal{I}}}
    \sum_{i\in\mathcal{I}}\omega_{y_i}\log p_i^m[y_i],
    \label{eq:rgrpl_student_loss}
\end{align}
where
\(Z_{\mathcal{I}}=\sum_{i\in\mathcal{I}}\omega_{y_i}\).
The distillation term is
\begin{equation}
    \mathcal{L}_{\mathrm{KD}}
    =\frac{1}{|\mathcal{I}|}
    \sum_{m\in\mathcal{M}}\sum_{i\in\mathcal{I}}
    D_{\mathrm{KL}}\!\left(
    p_i^{\mathrm{SDT},(\tau_s)}
    \,\Vert\,p_i^{m,(\tau_s)}\right),
    \label{eq:rgrpl_kd_loss}
\end{equation}
where \(\tau_s\) is the distillation temperature. The proposed method
preserves the SDT encoder, enhanced-feature fusion, self-distillation, and
final classifier, while introducing prototype-based geometric supervision
over the pure modality representations.

\subsection{Learnable Full-Dimensional Prototype Learning}
\label{subsec:rgrpl}

\subsubsection{Overview}

The proposed framework augments SDT with a shared bank of learnable emotion
prototypes in a complete \(D\)-dimensional space. Three modality-specific
projectors map pure text, audio, and visual features into a Poincar\'e ball,
where \(D=16\). Each class prototype combines an angular Emotion-Wheel
reference with a learnable residual orthogonal to the reference plane.
Prototype cross-entropy supplies discriminative supervision. Anchor
regularization limits prototype drift, while an angular relational loss
transfers the neighborhood structure of the Emotion Wheel to the learned
hyperbolic space.

The prototype branch provides an auxiliary geometric learning signal, while
the SDT classifier remains responsible for the final categorical prediction.

\subsubsection{Poincar\'e Projection of Modality Features}

We employ the curvature-\(-1\) Poincar\'e ball
\begin{equation}
    \mathbb{D}^{D}
    =\{z\in\mathbb{R}^{D}:\lVert z\rVert_2<1\}.
    \label{eq:rgrpl_poincare_ball}
\end{equation}
For modality \(m\), an affine map projects a pure feature into the tangent
space at the origin:
\begin{equation}
    v_i^m=W_m h_i^m+b_m,
    \qquad v_i^m\in T_0\mathbb{D}^{D}\simeq\mathbb{R}^{D}.
    \label{eq:rgrpl_feature_tangent}
\end{equation}
The exponential map produces the hyperbolic utterance representation
\begin{equation}
    z_i^m=\operatorname{Exp}_0(v_i^m)
    =\tanh(\lVert v_i^m\rVert_2)
    \frac{v_i^m}{\lVert v_i^m\rVert_2}
    \in\mathbb{D}^{D},
    \label{eq:rgrpl_feature_expmap}
\end{equation}
with the value at \(v_i^m=0\) defined by continuity. Sample radii are not
fixed, allowing the model to use both angular and radial information. The
geodesic distance between \(x,y\in\mathbb{D}^{D}\) is
\begin{equation}
    d_{\mathbb{D}}(x,y)
    =\operatorname{arcosh}\!\left(
    1+2\frac{\lVert x-y\rVert_2^2}
    {(1-\lVert x\rVert_2^2)(1-\lVert y\rVert_2^2)}
    \right).
    \label{eq:rgrpl_poincare_distance}
\end{equation}

\subsubsection{Angular Emotion References}

Each emotion class \(c\) is assigned an angular position \(\theta_c\) on the
Emotion Wheel. Its two-dimensional reference is
\begin{equation}
    q_c=r_0
    \begin{bmatrix}
        \cos\theta_c & \sin\theta_c
    \end{bmatrix}^{\top}
    \in\mathbb{R}^{2},
    \label{eq:rgrpl_angular_reference}
\end{equation}
where \(r_0>0\) determines the reference scale in the tangent space. Rather
than padding \(q_c\) with zeros, we introduce a learnable orthonormal map
\begin{equation}
    U\in\mathbb{R}^{D\times 2},
    \qquad U^\top U=I_2,
    \label{eq:rgrpl_psychological_subspace}
\end{equation}
which embeds the angular reference into a two-dimensional subspace of
\(T_0\mathbb{D}^{D}\). The matrix \(U\) is learned jointly with the network
subject to the orthonormality constraint. Consequently, the reference plane
can orient itself according to the multimodal representations, while the
angular inner products and distances among \(\{q_c\}_{c=1}^{C}\) are
preserved by the isometry. The embedded tangent-space anchor is
\begin{equation}
    \bar u_c=Uq_c.
    \label{eq:rgrpl_tangent_anchor}
\end{equation}

\subsubsection{Learnable Full-Dimensional Emotion Prototypes}

For each class, we introduce a learnable vector
\(r_c\in\mathbb{R}^{D}\). Its component within the angular reference subspace
is removed:
\begin{equation}
    r_c^\perp=(I_D-UU^\top)r_c.
    \label{eq:rgrpl_orthogonal_residual}
\end{equation}
The resulting prototype is obtained by mapping the combined tangent vector to
the Poincar\'e ball:
\begin{equation}
    \boxed{
    p_c=\operatorname{Exp}_0\!\left(Uq_c+\alpha r_c^\perp\right)}
    \in\mathbb{D}^{D}.
    \label{eq:rgrpl_learned_prototype}
\end{equation}
where \(\alpha\geq 0\) controls the degree of data-driven adaptation. The
term \(\bar u_c\) preserves the angular Emotion-Wheel reference, while
\(r_c^\perp\) provides \(D-2\) directions for learning discriminative class
structure that is not specified by the two-dimensional angular prior. A single
prototype bank \(\mathcal{P}=\{p_c\}_{c=1}^{C}\) is shared by all three
modalities, preventing independently drifting modality-specific class
geometries.

The construction separates two complementary roles. Learning \(U\) changes
the orientation of the angular reference plane but does not turn its
two-dimensional structure into \(D\) independent directions. The residuals
\(\{r_c^\perp\}\) provide the additional degrees of freedom that allow the
class prototypes to occupy the full ambient space. The orthogonality
constraint prevents those residuals from overwriting the angular coordinates
within \(\operatorname{span}(U)\).

\subsubsection{Anchor-Preserving Regularization}

Unconstrained residuals could dominate the angular component and cause
the prototypes to lose their intended interpretation. We therefore penalize
their squared norm:
\begin{equation}
    \mathcal{L}_{\mathrm{anchor}}
    =\frac{1}{C}\sum_{c=1}^{C}
    \lVert r_c^\perp\rVert_2^2.
    \label{eq:rgrpl_anchor_loss}
\end{equation}
This term regulates, rather than eliminates, adaptation. An excessively large
weight would shrink every residual toward zero and collapse the prototype
bank back to an embedded two-dimensional angular prior.

\subsubsection{Prototype Classification Objective}

For utterance \(i\), modality \(m\), and class \(c\), the prototype score is
the negative Poincar\'e distance:
\begin{equation}
    s_{i,c}^{m}
    =-\frac{d_{\mathbb{D}}(z_i^m,p_c)}{\tau_p},
    \label{eq:rgrpl_prototype_score}
\end{equation}
where \(\tau_p>0\) is the prototype temperature. Defining
\begin{equation}
    \rho_{i,c}^{m}
    =\frac{\exp(s_{i,c}^{m})}
    {\sum_{k=1}^{C}\exp(s_{i,k}^{m})},
    \label{eq:rgrpl_prototype_probability}
\end{equation}
the prototype cross-entropy is
\begin{equation}
    \mathcal{L}_{\mathrm{ProtoCE}}
    =-\frac{1}{|\mathcal{M}|\,|\mathcal{I}|}
    \sum_{m\in\mathcal{M}}\sum_{i\in\mathcal{I}}
    \log\rho_{i,y_i}^{m}.
    \label{eq:rgrpl_prototype_ce}
\end{equation}
This objective jointly updates the modality projectors, the relevant SDT
encoder features, and the learnable prototype residuals. Sharing the
prototypes across modalities also supplies a common class-level semantic
reference without requiring direct equality between heterogeneous modality
features.

\subsubsection{Angular Relational Learning in the Full-Dimensional Space}

Prototype classification alone does not specify how emotion classes should be
organized relative to one another. We therefore transfer the angular geometry
used by CPCC on the Emotion Wheel to the learned prototype bank.
For classes \(i\neq j\), the normalized shortest angular distance is
\begin{equation}
    D_{ij}^{W}
    =\frac{\arccos\!\left(\cos(\theta_i-\theta_j)\right)}{\pi},
    \qquad D_{ij}^{W}\in[0,1].
    \label{eq:rgrpl_wheel_distance}
\end{equation}
For every anchor class \(i\), these distances define a target relational
distribution over all other classes:
\begin{equation}
    Q_{ij}
    =\frac{\exp(-D_{ij}^{W}/\tau_W)}
    {\sum_{k\neq i}\exp(-D_{ik}^{W}/\tau_W)},
    \qquad j\neq i,
    \label{eq:rgrpl_wheel_neighbors}
\end{equation}
where \(\tau_W>0\) controls the concentration of the angular neighborhood.

The corresponding relations in the learned space are derived from the
Poincar\'e distances between full-dimensional prototypes:
\begin{align}
    D_{ij}^{P}
    &=d_{\mathbb{D}}(p_i,p_j),
    \label{eq:rgrpl_prototype_distance}\\
    P_{ij}
    &=\frac{\exp(-D_{ij}^{P}/\tau_P)}
    {\sum_{k\neq i}\exp(-D_{ik}^{P}/\tau_P)},
    \qquad j\neq i,
    \label{eq:rgrpl_prototype_neighbors}
\end{align}
where \(\tau_P>0\) controls the concentration of the learned prototype
neighborhood. Self-relations are masked in both \(Q\) and \(P\). Their
discrepancy is minimized through
\begin{equation}
    \boxed{
    \mathcal{L}_{\mathrm{relation}}
    =\frac{1}{C}\sum_{i=1}^{C}
    D_{\mathrm{KL}}\!\left(
    Q_{i,\neg i}\,\Vert\,P_{i,\neg i}\right)}.
    \label{eq:rgrpl_relation_loss}
\end{equation}
This objective preserves relative angular neighborhoods. It neither confines
the prototypes to a two-dimensional circle nor forces hyperbolic distances to
equal angular distances. Its influence is controlled by \(\lambda_r\); an
excessive weight can keep the learned geometry unnecessarily close to the
original wheel and leave the additional residual directions underused.

\subsubsection{Overall Learning Objective}

The complete learning objective is
\begin{equation}
    \boxed{\begin{aligned}
    \mathcal{L}_{\mathrm{total}}
    ={}&\mathcal{L}_{\mathrm{SDT}}
    +\lambda_p\mathcal{L}_{\mathrm{ProtoCE}}\\
    &+\lambda_r\mathcal{L}_{\mathrm{relation}}
    +\lambda_a\mathcal{L}_{\mathrm{anchor}}
    \end{aligned}}.
    \label{eq:rgrpl_objective}
\end{equation}
Here, \(\lambda_p\), \(\lambda_r\), and \(\lambda_a\) balance class-level
prototype discrimination, angular relational preservation, and
anchor-preserving regularization, respectively. The resulting objective
jointly shapes modality representations and class prototypes while retaining
the SDT categorical decision function
\begin{equation}
    \hat y_i
    =\underset{c\in\mathcal{Y}}{\arg\max}\;
    p_i^{\mathrm{SDT}}[c],
    \label{eq:rgrpl_prediction}
\end{equation}
for the final emotion prediction.

\subsubsection{Geometric Interpretation}

The angular prior remains two-dimensional even though its embedding matrix
\(U\) is learnable. The \(D\)-dimensional Poincar\'e ball is the learned
representation space, not a collection of \(D\) independent angular axes.
The plane \(\operatorname{span}(U)\) carries the original wheel references,
whereas its orthogonal complement supplies class-specific directions through
\(r_c^\perp\). Consequently, the high-dimensional capacity arises from the
learnable residuals rather than from rotating the reference plane alone. The
angular relational objective constrains only relative neighborhoods, leaving
the residual space available to encode discriminative structure supported by
the multimodal data.

\section{Experimental Setup}
\label{sec:experimental_setup}

\section{Experimental Results}
\label{sec:experimental_results}

\section{Conclusion}
\label{sec:conclusion}

\end{document}
